\documentclass[11pt]{amsart}

\usepackage[T1]{fontenc}
\usepackage{lmodern}
\usepackage[margin=1in]{geometry}
\usepackage{amsmath,amssymb,mathtools}
\usepackage{microtype}
\usepackage[hidelinks]{hyperref}
\usepackage{url}

\allowdisplaybreaks[2]

\newtheorem{theorem}{Theorem}[section]
\newtheorem{proposition}[theorem]{Proposition}
\newtheorem{lemma}[theorem]{Lemma}
\newtheorem{corollary}[theorem]{Corollary}
\theoremstyle{remark}
\newtheorem{remark}[theorem]{Remark}

\DeclareMathOperator{\tr}{tr}
\newcommand{\R}{\mathbb R}
\newcommand{\D}{\mathcal D}
\newcommand{\F}{\mathrm F}

\title[Two-certificate trace--energy deduction]
{A two-certificate trace--energy deduction for simple zeros of the
Riemann zeta function}
\author{Yuhang Shi}
\email{yuhangshi888@gmail.com}
\date{August 14, 2026}

\begin{document}
\begin{abstract}
Let $N(T,2T)$ count the nontrivial zeros of the Riemann zeta function with
ordinates in $(T,2T]$, with multiplicity, and let $N_0^s(T,2T)$ count the
simple zeros on the critical line in that interval.  We refine the
finite-block deduction used in recent stability improvements of the
rank--trace method.  The pressure term in a local gap certificate can be
absorbed by a supporting plane obtained from two certified local gap
inequalities and the finite-dimensional trace--energy envelope.  Applied to
the certified seven- and nine-point inequalities on the common window of
\texttt{trmdy/zeta-simple-zeros-673137}, the refinement gives the
research-draft candidate
\[
 \liminf_{T\to\infty}\frac{N_0^s(T,2T)}{N(T,2T)}
 \ge 0.6733169771424713\ldots>0.673316977.
\]
No new large numerical search is needed: the only new input is a
finite-dimensional inequality proved here.  The analytic interface and the
machine-certified local inequality are imported at their explicitly stated
proof boundary.
\end{abstract}
\maketitle

\section{Statement and proof boundary}

Put
\[
 S=N_0^s(T,2T),\qquad N=N(T,2T).
\]
We use the arbitrary-window Gabor and stability interface developed in
\cite{Claude2026,Ainta2026} and instantiated with a perturbed admissible
window in \cite{Trmdy2026}.  In the notation of those sources, it gives
\begin{equation}\label{eq:interface}
 S\ge H_{\rm cert}N+\D(M)-o(N),
 \qquad
 H_{\rm cert}=\frac{3362285207}{5000000000},
\end{equation}
where $M$ is the Gram matrix of the retained simple-zero atoms and
\[
 \D(G)=\tr\Psi(G),\qquad
 \Psi(t)=
 \begin{cases}
  (t-1)^2,&0\le t\le2,\\
  2t-3,&t\ge2.
 \end{cases}
\]
For every fixed $R_0>0$, the entries of a fixed-size principal Gram block
satisfy, uniformly when all its normalized separations are at most $R_0$,
\begin{equation}\label{eq:gram-limit}
 G_{ij}=k_v(y_i-y_j)+o(1),
\end{equation}
where $k_v$ is the normalized overlap kernel of the chosen window.  We also
use convex spectral pinching and the shifted consecutive-block
decomposition of \cite{Ainta2026,Trmdy2026}.

The finite inputs are two certified inequalities from \cite{Trmdy2026}.
Let $w=k_v^2$.  For $q\in\{6,8\}$, let $g_1,\ldots,g_q\ge0$, put $y_0=0$
and $y_j=g_1+\cdots+g_j$, and define
\[
 (\epsilon_q,p_q)=
 \begin{cases}
  (891/200000,1/2736),&q=6,\\
  (15211/2500000,1/2500),&q=8.
 \end{cases}
\]
There are nonnegative rational numbers $a_{ij}^{(q)}$ such that
\begin{align}
 p_q\sum_{r=1}^qg_r
 +\sum_{0\le i<j\le q}a_{ij}^{(q)}w(y_j-y_i)
 &\ge \epsilon_q,
 \label{eq:local}\\
 \sum_{i=0}^{q-r}a_{i,i+r}^{(q)}&=2
 \qquad(1\le r\le q).
 \label{eq:capacity}
\end{align}
The exact $21$ and $36$ weights and their exhaustive interval certificates
are contained in the repository at commit
\texttt{1610b97b7895ff34982260f8dcaf04a0f7b82cf7}.  Our result is conditional
only in the sense that \eqref{eq:interface}, \eqref{eq:gram-limit}, and the
machine certificate \eqref{eq:local} are imported rather than reproved here.

\begin{theorem}\label{thm:main}
Subject to the imported interfaces and certificate just specified,
\[
 \liminf_{T\to\infty}\frac{N_0^s(T,2T)}{N(T,2T)}
 \ge B,
\]
where, with
\begin{align*}
 A_7&=\frac{189783}{200000},
 \qquad A_9=\frac{3209521}{2500000},\\
 R&=2\sqrt{\frac{349837789}{273750000}}-1
       +\frac{3209521}{547500000},\\
 u&=\frac{R-A_7}{A_9-A_7},\\
 B&=
 \frac{219H_{\rm cert}
 -(1-u)\,213\cdot6/2736-u\,211\cdot8/2500}{219-R},
\end{align*}
one has
\[
 B=0.6733169771424713134\ldots
 >\frac{673316977}{10^9}.
\]
\end{theorem}

The preceding public assembly of the nine-point certificate gives
$0.6733127422722459981\ldots$ \cite{Trmdy2026}; thus the gain in proportion
is approximately $4.2349\times10^{-6}$.  All such August 2026 follow-up
values are very recent and should be regarded as research-draft candidates
pending independent mathematical review.

\section{The scaled pressure lemma}

For $m\ge3$, define
\begin{equation}\label{eq:phi}
 \Phi_m(E)=
 \begin{cases}
  E,&0\le E\le m/(m-1),\\[2mm]
  2\sqrt{\dfrac{m-1}{m}E}-1+\dfrac Em,
    &E\ge m/(m-1).
 \end{cases}
\end{equation}
The function is increasing and concave.  Indeed, it is linear on the first
branch, has negative second derivative on the second branch, and the two
one-sided derivatives agree and equal $1$ at $E=m/(m-1)$.

\begin{lemma}[Scaled finite-dimensional envelope]\label{lem:scaled}
Let $G$ be an $m\times m$ positive-semidefinite matrix with unit diagonal.
If its eigenvalues are $\lambda_1,\ldots,\lambda_m$, put
\[
 E=\sum_{i=1}^m(\lambda_i-1)^2,
 \qquad D=\tr\Psi(G).
\]
Let $A\ge m/(m-1)$ satisfy $\Phi_m(A)<2$, set
\[
 R=\Phi_m(A),\qquad \eta=\frac RA,
\]
and let $P\ge0$.  If $E+P\ge A$, then
\begin{equation}\label{eq:scaled}
 D+\eta P\ge R.
\end{equation}
\end{lemma}

\begin{proof}
Write $x_i=\lambda_i-1$.  Then $x_i\ge-1$, $\sum_i x_i=0$, and
$E=\sum_i x_i^2$.  Let
\[
 L=\{i:x_i>1\},\quad q=|L|,
 \quad X=\sum_{i\in L}x_i,
 \quad Q=\sum_{i\in L}x_i^2.
\]
If $q>0$, then $q<m$, because $\sum_i x_i=0$.
Since the quadratic branch $x_i^2$ is replaced by $2x_i-1$ precisely on
$L$,
\begin{equation}\label{eq:Didentity}
 D=E+2X-q-Q.
\end{equation}

If $q\ge2$, Cauchy--Schwarz on the $m-q$ complementary coordinates gives
$E-Q\ge X^2/(m-q)$.  Since $X>q$,
\begin{equation}\label{eq:qge2}
 D\ge2X-q+\frac{X^2}{m-q}
 >q+\frac{q^2}{m-q}=\frac{qm}{m-q}>2>R.
\end{equation}

Suppose next that $q=1$, and write $x_i=r>1$ for the unique large
coordinate.  Cauchy--Schwarz on the remaining $m-1$ coordinates gives
\[
 E\ge r^2+\frac{r^2}{m-1},
 \qquad
 r\le\sqrt{\frac{m-1}{m}E}.
\]
In particular, the right endpoint in the latter inequality is greater than
$1$.
By \eqref{eq:Didentity},
\[
 D=E+2r-1-r^2.
\]
The function $2r-r^2$ decreases for $r\ge1$, and therefore
\begin{equation}\label{eq:one-large}
 D\ge\Phi_m(E).
\end{equation}
If $q=0$, then simply $D=E$; moreover $q=0$ necessarily holds when
$E<m/(m-1)$.

The concavity of $\Phi_m$ and $\Phi_m(0)=0$ imply the chord inequality
\begin{equation}\label{eq:chord}
 \Phi_m(E)\ge\frac EA\Phi_m(A)=\eta E
 \qquad(0\le E\le A).
\end{equation}
Also $R\le A$, since on the second branch
\[
 A-\Phi_m(A)
 =\left(\sqrt{\frac{m-1}{m}A}-1\right)^2\ge0.
\]
If $E<A$, then \eqref{eq:qge2}, \eqref{eq:one-large}, the $q=0$ case, and
\eqref{eq:chord} give either $D>R$ or $D\ge\eta E$.  In the latter case
$D+\eta P\ge\eta(E+P)\ge R$.  If $E\ge A$, the same case analysis gives
$D\ge R$: this is immediate when $q\ge2$; when $q=1$ it follows from
\eqref{eq:one-large} and monotonicity of $\Phi_m$; and when $q=0$ it follows
from $D=E\ge A\ge R$.  This proves \eqref{eq:scaled}.
\end{proof}

The next form retains two independent local constraints before paying their
global pressure costs.

\begin{lemma}[Two-certificate supporting plane]\label{lem:two}
Let $G$, $E$, and $D$ be as in Lemma~\ref{lem:scaled}.  Suppose that
\[
 E+p_7L_7\ge A_7,\qquad E+p_9L_9\ge A_9,
 \qquad L_7\ge\frac34L_9,
\]
where $L_7,L_9\ge0$, $0<A_7<m/(m-1)<A_9$, and
$R=\Phi_m(A_9)<2$.  Put
\[
 u=\frac{R-A_7}{A_9-A_7},\qquad
 \beta=(1-u)p_7,qquad \gamma=up_9.
\]
If $0\le u\le1$, $A_7/p_7\ge3A_9/(4p_9)$, and
\begin{equation}\label{eq:slope-condition}
 \frac{3\beta}{4p_9}+\frac{\gamma}{p_9}\le1,
\end{equation}
then
\begin{equation}\label{eq:two-plane}
 D+\beta L_7+\gamma L_9\ge R.
\end{equation}
\end{lemma}

\begin{proof}
The case $D>R$ is immediate.  In every remaining case of the proof of
Lemma~\ref{lem:scaled}, $D\ge\Phi_m(E)$.  The two hypotheses and
$L_7\ge3L_9/4$ give
\begin{align*}
 L_9&\ge\max\!\left(\frac{A_9-E}{p_9},0\right),\\
 L_7&\ge
 \max\!\left(\frac{A_7-E}{p_7},
       \frac34\frac{A_9-E}{p_9},0\right).
\end{align*}
It is thus enough to minimize the resulting lower bound as a function of
$E\ge0$.

On $[0,A_7]$ the first branch of \eqref{eq:phi} applies.  The lower bound is
piecewise affine.  At $E=0$ the active lower bound for $L_7$ is
$(A_7-E)/p_7$ by the extra numerical hypothesis; when the active affine
branch changes, the slope remains nonnegative by
\eqref{eq:slope-condition} and by
\[
 \frac{\beta}{p_7}+\frac{\gamma}{p_9}=1.
\]
Its value at $E=0$ is $R$.  On $[A_7,A_9]$ it is
\[
 \Phi_m(E)+\left(\frac{3\beta}{4}+\gamma\right)
 \frac{A_9-E}{p_9}.
\]
More precisely, the first interval starts at the value $R$ and has
nonnegative slope, so its value at $A_7$ is at least $R$.  The latter
function is concave, has that same value at $A_7$, and equals $R$ at $A_9$;
hence it is at least $R$ throughout.  For $E\ge A_9$, monotonicity of
$\Phi_m$ gives $\Phi_m(E)\ge R$.  This proves \eqref{eq:two-plane}.
\end{proof}

\begin{remark}
The slope $\eta$ is optimal for a conclusion of the form
$D+cP\ge\Phi_m(A)$ based only on $E+P\ge A$: at $E=0$, $P=A$, one must have
$cA\ge\Phi_m(A)$.  Thus the improvement below is the full gain available
from rescaling this single pressure term at a fixed block length.
\end{remark}

\section{Application to the seven- and nine-point certificates}

Let $S^\circ$ denote the number of retained central simple zeros.  The
endpoint trimming in the imported interface gives
\begin{equation}\label{eq:retained}
 S^\circ=S-o(N).
\end{equation}
Fix $m=219$ and partition these retained zeros into consecutive $m$-point
blocks.  For a block
\[
 y_1<\cdots<y_m,
\]
let $W_{q,B}$ be the sum of the spans of all consecutive $(q+1)$-point
windows fully contained in the block:
\[
 W_{q,B}=\sum_{j=1}^{m-q}(y_{j+q}-y_j)\qquad(q=6,8).
\]
Summing \eqref{eq:local} over these $m-q$ windows and using
\eqref{eq:capacity}, every pair term receives total coefficient at most
$2$.  Hence, with
\[
 E_B=2\sum_{1\le i<j\le m}w(y_j-y_i),
\]
we have the two simultaneous estimates
\begin{equation}\label{eq:block-energy}
 E_B+p_qW_{q,B}\ge A_q:=\epsilon_q(m-q)
 \qquad(q=6,8).
\end{equation}
Moreover, coefficientwise in the individual gaps,
\begin{equation}\label{eq:span-comparison}
 W_{6,B}\ge\frac34W_{8,B}.
\end{equation}
Indeed, an interior gap occurs six times in $W_{6,B}$ and eight times in
$W_{8,B}$, while at either endpoint the two multiplicities increase together
from one to six.

For $m=219$,
\[
 A_7=\frac{189783}{200000}<\frac{219}{218}
 <A_9=\frac{3209521}{2500000},
\]
and
\begin{equation}\label{eq:R}
 R=\Phi_{219}(A_9)
 =2\sqrt{\frac{349837789}{273750000}}-1
   +\frac{3209521}{547500000}
 =1.2667878440823898\ldots<2.
\end{equation}
Put $u=(R-A_7)/(A_9-A_7)$,
$\beta=(1-u)/2736$, and $\gamma=u/2500$.
Then $0<u<1$, and
\[
 \frac{\beta}{1/2736}+\frac{\gamma}{1/2500}=1.
\]
Also $3\beta/(4p_9)+\gamma/p_9<1$: as a function of $u$ its slope is
$1-3p_7/(4p_9)=287/912>0$, and it equals $1$ at $u=1$, whereas $u<1$.
Finally,
\[
 \frac{A_7}{p_7}-\frac{3A_9}{4p_9}
 =\frac{18909069}{100000}>0.
\]

We claim that every full block satisfies
\begin{equation}\label{eq:block-defect}
 \D(G_B)+\beta W_{6,B}+\gamma W_{8,B}\ge R-o(1),
\end{equation}
uniformly in the block.  If
$\beta W_{6,B}+\gamma W_{8,B}\ge R$, this is automatic.  Otherwise
$W_{8,B}<R/\gamma$, since $\gamma>0$, so $W_{8,B}$ is uniformly bounded.
Because every gap in the block occurs in at least one of its nine-point
windows,
\[
 y_m-y_1\le W_{8,B};
\]
thus the whole block has uniformly bounded span.  The compact-uniform
asymptotic \eqref{eq:gram-limit}, with fixed $m$, gives
\[
 2\sum_{i<j}|(G_B)_{ij}|^2=E_B+o(1)
\]
uniformly.  The diagonal entries are $1+o(1)$ in the imported interface.
Let $\widetilde G_B$ be the correlation normalization of $G_B$.  It is
positive semidefinite with unit diagonal and, because $m$ is fixed, replacing
$G_B$ by $\widetilde G_B$ changes its finitely many entries, its energy, and
$\D$ by $o(1)$ uniformly.  For $\widetilde G_B$ the energy is exactly
$\operatorname{tr}(\widetilde G_B-I)^2
=2\sum_{i<j}|(\widetilde G_B)_{ij}|^2$.
Lemma~\ref{lem:two}, continuity of $\Phi_{219}$, and
\eqref{eq:block-energy}--\eqref{eq:span-comparison} then prove
\eqref{eq:block-defect}.

It remains to average \eqref{eq:block-defect}.  Use the $m$ shifted
partitions into consecutive $m$-point blocks, retaining an endpoint block of
at most $2(m-1)$ points for each shift.  Convex spectral pinching gives
\[
 \D(M)\ge\sum_B\D(G_B)
\]
for every shift.  A fixed consecutive $(q+1)$-point window is fully contained
in a block for at most $m-q$ of the $m$ shifts, with equality away from the
two finite endpoints.  Endpoint omissions only decrease the pressure sums,
whereas the number of full blocks, averaged over the shifts, is
$S^\circ/m+O_m(1)$.  Moreover,
\begin{equation}\label{eq:window-count}
 \sum_j(y_{j+q}-y_j)
 \le q(y_{S^\circ}-y_1)\qquad(q=6,8),
\end{equation}
because every gap occurs in at most $q$ consecutive $(q+1)$-point windows.
The normalized total span is at most $N+o(N)$ by the
Riemann--von Mangoldt formula.  The blockwise $o(1)$ in
\eqref{eq:block-defect} is uniform and the number of blocks is $O(S)$.
Consequently
\begin{equation}\label{eq:global-defect}
 \D(M)\ge
 \frac{R}{m}S^\circ
 -\frac{\beta\,6(m-6)+\gamma\,8(m-8)}mN-o(N).
\end{equation}
Using \eqref{eq:retained} replaces $S^\circ$ by $S$ at a cost $o(N)$.

Substituting \eqref{eq:global-defect} into \eqref{eq:interface} and
rearranging gives
\begin{equation}\label{eq:bound}
 \frac SN\ge
 \frac{mH_{\rm cert}-\beta\,6(m-6)-\gamma\,8(m-8)}{m-R}-o(1).
\end{equation}
At $m=219$, this is the constant $B$ in Theorem~\ref{thm:main}.

For completeness, the strict decimal comparison can be checked using
rational arithmetic and one squaring.  From \eqref{eq:R},
\[
 R>R_0:=\frac{6333939}{5000000},
\]
because
\[
 \frac{349837789}{273750000}
 -\left(\frac{R_0+1-A_9/219}{2}\right)^2
 =\frac{239006467199}{4796100000000000000}>0.
\]
Since
\[
 \beta\,6(m-6)+\gamma\,8(m-8)
 =\frac{19769000}{31814873}R-\frac{389820601}{3181487300},
\]
the right-hand side of \eqref{eq:bound} is increasing in $R$: its derivative
has the sign of
\[
 \frac{1799014260305932709}{159074365000000000}>0.
\]
Replacing $R$ by $R_0$ therefore gives the rational comparison
\[
 \frac{23320853620214932709}{34635772470125253000}
 -\frac{673316977}{10^9}
 =\frac{4570374547679819}{34635772470125253000000000}>0.
\]
This proves the final strict inequality without relying on rounded decimal
arithmetic.

\section{Relation to the preceding deduction}

The preceding public finite-block argument uses one local certificate at a
time.  Lemma~\ref{lem:two} retains the simultaneous consequences of the
certified seven- and nine-point inequalities before their pressure terms are
paid globally.  The parameter $u$ is the barycentric coordinate of $R$
between $A_7$ and $A_9$, so the resulting pressure tax is a weighted
combination of two independently certified taxes.  No window, local weight,
or interval certificate has been changed.

For comparison, the one-certificate implication is
\[
 E+P\ge A\quad\Longrightarrow\quad
 D+P\ge\Phi_m(A).
\]
Lemma~\ref{lem:scaled} replaces the coefficient of $P$ by
$\eta=\Phi_m(A)/A<1$.  In the shifted-block average this decreases the
global pressure tax while preserving the same block defect $R$.  No window,
local weights, or interval certificate has been changed.

The Schur--Jensen remainder in \cite{Shi2026} is not added separately here.
Its scalar function $g$ satisfies $\Psi(t)=g(t)+1$.  If
$d_j=M_{jj}\le1$ and
\(
 \mathcal S_g=\tr g(M)-\sum_jg(d_j),
\)
then the exact identity
\[
 \tr\Psi(M)-\mathcal S_g
 =\sum_j\bigl(1+g(d_j)\bigr)
 =\sum_j(1-d_j)^2\ge0
\]
holds.  In particular, the two quantities agree for a unit-diagonal Gram
matrix.
Adding both terms would therefore double count the same stability slack.
The earlier work motivated retaining this information, but the present gain
comes specifically from a sharper use of the already retained
$\tr\Psi$ defect.

\section*{Acknowledgements}

The author thanks the contributors of the certified local inequalities and
window in \cite{Trmdy2026}.  The finite-dimensional envelope and
window-in-frame counting used there trace to \cite{Tawan2026}.

\begin{thebibliography}{99}

\bibitem{Claude2026}
Claude,
\emph{More than two thirds of the zeros of the Riemann zeta function lie on
the critical line}, preprint, August 10, 2026; paper and Lean~4 artifact,
\url{https://github.com/anthropics/zeta-23-lean}.

\bibitem{Ainta2026}
\texttt{ainta/zeta-simple-zeros},
\emph{More than $67.3\%$ of the zeros of the Riemann zeta function are
simple and lie on the critical line}, research repository, 2026,
\url{https://github.com/ainta/zeta-simple-zeros}.

\bibitem{Tawan2026}
\path{tawanerguo-cn/zeta-simple-zeros},
\emph{A certified unconditional $67.3192911473\%$ lower bound for simple
zeros of the Riemann zeta function}, research repository, 2026,
\url{https://github.com/tawanerguo-cn/zeta-simple-zeros}.

\bibitem{Trmdy2026}
\texttt{trmdy/zeta-simple-zeros-673137},
\emph{A $67.3312742272\%$ lower-bound candidate for simple zeros of zeta},
research repository, commit
\texttt{1610b97b7895ff34982260f8dcaf04a0f7b82cf7}, 2026,
\url{https://github.com/trmdy/zeta-simple-zeros-673137}.

\bibitem{Shi2026}
Y.~Shi,
\emph{A Schur--Jensen Gain in the Critical-Line Zero Problem},
Zenodo, 2026,
\url{https://doi.org/10.5281/zenodo.21903013}.

\bibitem{Montgomery1973}
H.~L.~Montgomery,
\emph{The pair correlation of zeros of the zeta function}, in
\emph{Analytic Number Theory}, Proc. Sympos. Pure Math. \textbf{24},
Amer. Math. Soc., 1973, 181--193.

\end{thebibliography}

\end{document}
